\documentclass[conference]{IEEEtran}
\IEEEoverridecommandlockouts

\usepackage{cite}
\usepackage{amsmath,amssymb,amsfonts}
\usepackage{graphicx}
\usepackage{textcomp}
\usepackage{xcolor}
\usepackage{tikz}
\usetikzlibrary{arrows.meta}
\usepackage{booktabs}
\usepackage{multirow}
\usepackage{placeins}
\usepackage{balance}
\usepackage{url}
\usepackage{array}
\usepackage{tabularx}
\usepackage{algorithm}
\usepackage{algpseudocode}
\usepackage[colorlinks=true,linkcolor=blue,citecolor=blue,urlcolor=blue]{hyperref}

\begin{document}

\title{Anonymity Preserving, Decoupled Distributed\\Authentication for Multihop IoT Networks}

\author{
  \IEEEauthorblockN{Sreeparna Das, Anup Kulkarni,  Manas Khatua}
  \IEEEauthorblockA{\textit{Computer Science and Engineering}\\
  \textit{Indian Institute of Technology Guwahati}\\
  \{d.sreeparna, anup.kulkarni, manaskhatua\}@iitg.ac.in}
  
}

\maketitle

% -------------------------------------------------------
\begin{abstract}
The Internet of Things (IoT) demands lightweight and secure authentication
for resource constrained multihop networks, while also protecting the
privacy of individual devices.  Based on our analysis of the lightweight decoupled distributed
authentication scheme of Das~\cite{das2026comsnets}, we identify
two critical gaps. First, the real device identity is exposed on the open
channel in every authentication session, enabling adversarial cross-session
tracking. Second, packet loss in the session key establishment phase desynchronises
the device and server states, forcing full re-enrolment.  We propose a security improved and anonymity preserving enhanced version of Das~\cite{das2026comsnets},
addressing both identified gaps while preserving the core features of the
original scheme, including PUF-based secret derivation, hash-based accumulator
membership verification, decoupled authentication server selection, and the
lightweight hash-XOR computational structure.  Device anonymity is preserved through per-session
pseudonym rotation, where a rotating pseudonym replaces the real device identity
in all open-channel communications, preventing long-term tracking.  A
dual-state synchronization mechanism maintained at both the authentication
server and the gateway enables seamless recovery from single packet-loss
events without re-enrolment.  Security is validated formally (ProVerif) and informally.
COOJA/Contiki-NG simulations confirm that the proposed scheme achieves a
per-device energy of 52.21\,mJ and a communication cost of 1376\,bits,
reducing total energy by 42.8\,\% over LAAKA~\cite{ali2024laaka} and 34.6\,\% over Zhou~\cite{zhou2024auth}, while
maintaining the lowest overhead among all compared schemes.
\end{abstract}


\begin{IEEEkeywords}
Lightweight, Authentication, Anonymity, Pseudonym, Physically
Unclonable Function (PUF), Dual-State Storage, Proverif, COOJA.
\end{IEEEkeywords}

% -------------------------------------------------------
\section{Introduction}

The Internet of Things (IoT) connects millions of smart devices that collect
and transmit data to centralized platforms~\cite{zhou2021iot}
such as gateways, often through multihop paths~\cite{lajnef2022wsn}.
Industrial and outdoor deployments expose devices to a range of adversarial
threats, making authentication a prerequisite for any communication session.
Because IoT nodes are battery-powered and resource-constrained,
authentication schemes must be \emph{lightweight} as well as
secure~\cite{khatua2024puf}.

Centralized schemes route all authentication traffic through a gateway,
imposing excessive relay load on intermediate nodes in a multihop
topology~\cite{sheik2023auth}\cite{he2023lightweight}\cite{roy2025design}.  Distributed
schemes distribute this responsibility. However, a plain parent-node scheme
overloads whichever node a set of new devices happened to join under, while
all other capable nodes remain idle~\cite{rosa2023lightsae}.

Das~\cite{das2026comsnets} addressed this skewed-load problem by proposing a
lightweight decoupled distributed authentication scheme in which the gateway
pre-assigns any capable IoT node — parent or non-parent — as the
authentication server for newly joined devices.  By decoupling authentication
from network topology, the scheme distributes the authentication load evenly
across the network.  Authentication relies on PUF-based secret derivation, a
hash-based one-way accumulator for group membership verification, and
XOR-masked message exchanges, resulting in an efficient and lightweight
protocol.

Despite its efficiency, Das~\cite{das2026comsnets} leaves two security gaps
unaddressed.  First, the real device identity $ID_D$ is transmitted in
plaintext on the open channel in every authentication session, enabling an
adversary with radio access to track a device across sessions — a direct
violation of modern privacy requirements.  Second, both the authentication
server and the device maintain only a single session state; a packet loss
during session key delivery desynchronises their states, causing all
subsequent authentication attempts to fail until the device undergoes full
re-enrolment, which is an energy-costly operation for battery-powered
constrained nodes.

This paper makes the following key contributions.
\begin{enumerate}
  \item An improved scheme over Das~\cite{das2026comsnets} is proposed,
    adding per-session pseudonym rotation for device anonymity and a
    dual-state desynchronisation recovery mechanism.
  \item The security of the proposed scheme is validated through both
    informal analysis and formal verification using
    ProVerif.
  \item The efficiency of the proposed scheme is evaluated using both
    theoretical and simulation-based (COOJA/Contiki-NG) frameworks,
    emphasizing communication overhead, energy consumption, and CPU
    utilization in a multihop constrained setting.
\end{enumerate}

Section~II elaborates on related literature.  Section~III presents the
proposed scheme.  Furthermore, Section~IV contains the detailed security
analysis.  Finally, Section~V presents the performance analysis, with
Section~VI stating the conclusion.



% -------------------------------------------------------
\section{Literature Survey}

Existing schemes for IoT authentication address the problem from multiple
directions but each leaves gaps.  Distributed approaches reduce relay load but
remain topology-bound. Light-SAE~\cite{rosa2023lightsae}
authorise parent nodes as authentication servers, yet a popular parent is
overwhelmed by simultaneous join requests from many children.
Zheng~et~al.~\cite{zheng2022puf} designed a peer-to-peer PUF-based scheme
that lacks a multihop decoupled design. RAAF-MEC~\cite{alruwaili2025raaf}
deploy multiple PUF-equipped edge servers, but the absence of a genuine
multihop topology means certain servers still receive disproportionate load.
PUF-based schemes avoid long-term secret storage but trade in overhead.
Li~et~al.~\cite{li2020puf} and SAKMS~\cite{yang2024sakms} combine PUF
with session key exchange at the cost of ECC operations, while
He~et~al.~\cite{he2023lightweight} achieve lightweight anonymous authentication
using hash-and-XOR primitives for CoAP/OSCORE but rely on a centralized server.

Identity-hiding approaches introduce anonymity at prohibitive cost.
Al-Riyami and Paterson~\cite{alriyami2003cert} propose ECC-based certificateless
cryptography that is too heavy for constrained nodes, and blockchain-based
schemes~\cite{ahmed2023chains} achieve decentralisation but incur prohibitive
storage and consensus overhead.  Zhong~et~al.~\cite{zhong2025survey} present a
systematic taxonomy of lightweight anonymous IoT authentication, deconstructing
existing protocols into nine sub-frameworks (SF1-UA through SF9-FS) and two
general frameworks (GF1, GF2), and their analysis demonstrates that only GF2
simultaneously achieves user anonymity, forward secrecy, and desynchronization
resistance, a combination absent from all prior lightweight schemes.  The
proposed scheme follows the GF2 architectural pattern, where the dual-state pseudonym
mechanism corresponds to the SF6-UA sub-framework~\cite{zhong2025survey}, and
per-session key derivation from a PUF-assisted fresh nonce provides forward
secrecy analogous to GF2.  Unlike GF2, which uses a one-time hash chain
$hk = h(hk)$, the proposed scheme replaces the hash chain with PUF-derived
session randomness and augments membership verification with a hash-based
accumulator — mechanisms not modelled in the Zhong~et~al.\ taxonomy.
Table~\ref{tab:features} contrasts key features across representative schemes.
Existing lightweight schemes trade off at least one of these properties, namely
decoupled design, PUF, accumulator-based membership, authentication secret update
(ASU), anonymity, or desynchronization recovery, whereas the proposed scheme is
the first to satisfy all six simultaneously with low computational and communication overhead.

\begin{table}[!t]
  \centering
  \caption{Features Comparison of Authentication Schemes.}
  \label{tab:features}
  \scriptsize
  \setlength{\tabcolsep}{3.5pt}
  \begin{tabular}{lcccccccc}
    \toprule
    Scheme & Dec. & PUF & Acc. & ASU & Anon. & Desync & C.OH & M.OH \\
    \midrule
    RAAF-MEC~\cite{alruwaili2025raaf} & $\times$ & \checkmark & $\times$ & \checkmark & $\times$ & $\times$ & H & L \\
    SAKMS~\cite{yang2024sakms}     & \checkmark & $\times$ & $\times$ & $\times$ & $\times$ & $\times$ & H & L \\
    Li et al.~\cite{li2020puf}         & \checkmark & \checkmark & $\times$ & \checkmark & $\times$ & $\times$ & H & L \\
    Zheng et al.~\cite{zheng2022puf}      & \checkmark & \checkmark & $\times$ & \checkmark & $\times$ & $\times$ & M & L \\
    Das et al.~\cite{das2026comsnets}  & \checkmark & \checkmark & \checkmark & \checkmark & $\times$ & $\times$ & L & L \\
    Proposed                 & \checkmark & \checkmark & \checkmark & \checkmark & \checkmark & \checkmark & L & L \\
    \bottomrule
  \end{tabular}\\[4pt]
  \raggedright
  \scriptsize
  \textit{Dec.}: Decoupled design.\ \
  \textit{Acc.}: Hash-based accumulator.\ \
  \textit{ASU}: Auth.\ secret update.\\
  \textit{Anon.}: Device anonymity via pseudonym.\ \
  \textit{Desync}: Desync.\ recovery.\\
  \textit{C.OH}: Computation — L (Hash/XOR/$\le2$ AES), M ($>2$ AES), H (ECC/Pairings).\\
  \textit{M.OH}: Message exchanges — L (1--3), M (4--5), H ($>5$).
\end{table}

% -------------------------------------------------------
\section{Proposed Authentication Scheme}
\label{sec:proposed}

\subsection{System Model and Assumptions}
\label{sec:sysmodel}

The proposed scheme is designed to help nodes establish secure multihop
communications with the gateway.  To establish a secure communication, nodes
should be authenticated by a trusted authentication server.  A set of
existing IoT nodes inside the network is entrusted with the responsibility of
authentication by the gateway in a distributed manner.  The gateway also
shares a secret key with each of those responsible IoT
nodes/authentication servers.  The authentication server can be a parent or
a non-parent node—they may be \emph{decoupled}—which prevents newly joined
child nodes from overwhelming their parent with biased authentication load.

IoT nodes are assumed to be embedded with a PUF and are capable of performing
various cryptographic operations, such as encryption/decryption, XOR,
hashing, bitwise AND, and concatenation.

Nodes $D$ and $AS$ are assumed to be the newly joined node and its
authentication server, respectively.  $GW$ represents the gateway
symbolically.  $AS$ may be located at a singlehop or multihop distance from
$D$.  It is assumed that $GW$ has already selected the authentication server
set and has established a shared secret key with each member of the set.
Let $\mathcal{AS}' = \{AS_1, AS_2, \ldots, AS_n\}$ be the set of
authentication servers selected by $GW$, such that $AS \in \mathcal{AS}'$.
$GW$ has a shared secret key $K_{GW\text{-}AS}$ with $AS$.  It is also
assumed that $D$, whose distinct identifier is $ID_D$, possesses the address
of its authentication server $AS$ along with its unique identity $ID_{AS}$.
Table~\ref{tab:notations} elaborates symbols and functions incorporated in
the following scheme.

There are mainly three phases — node enrolment, node authentication, and node
session key exchange.  The following sections elaborate on these phases.

\begin{table}[!t]
  \centering
  \caption{Notation.}
  \label{tab:notations}
  \scriptsize
  \setlength{\tabcolsep}{4pt}
  \begin{tabular}{ll}
    \toprule
    Symbol & Description \\
    \midrule
    $D$, $AS$ & Newly joined device and its authentication server \\
    $GW$ & Gateway / root node \\
    $H(\cdot)$ & Collision-resistant hash function (SHA-256) \\
    $\mathsf{SE}(\cdot)$, $\mathsf{SD}(\cdot)$ & Symmetric encryption/decryption (AES-128) \\
    $\oplus$ & Bitwise XOR operation \\
    $\mathbin{\&}$ & Bitwise AND operation \\
    $c_D$, $c_{AS\text{-}D}$ & PUF challenges for $D$ (issued by $AS$) and $AS$ (by $D$) \\
    $m_D$ & Session-based random issued by $AS$ during enrolment \\
    $m_{\mathrm{curr}}$, $m_{\mathrm{old}}$ & Current and previous session-based randoms at $D$ and $AS$ \\
    $y_D$, $R_D$ & Group secret and PUF response of $D$ \\
    $Y_D^H$ & Hashed unique secret: $H(y_D)$ \\
    $T_\mathsf{Acc}$ & Accumulated group-membership value at $AS$ \\
    $\Phi_{AS\text{-}D}$ & PUF-response binding: $\mathsf{PUF}_{AS}(c_{AS\text{-}D}) \oplus R_D$ \\
    $\mathit{PID}_{\mathrm{curr}}$, $\mathit{PID}_{\mathrm{old}}$ & Current and previous pseudonyms of $D$ \\
    $K_{GW\text{-}D}$ & Session key between $GW$ and $D$ \\
    $K_{GW\text{-}AS}$ & Long-term shared key between $GW$ and $AS$ \\
    $ts_1$, $ts_2$, $ts_{\mathrm{auth}}$ & Timestamps for freshness verification \\
    \bottomrule
  \end{tabular}
\end{table}

\subsection{Node Enrolment Phase}

The main purpose of this phase is to enrol the unique secret of $D$ (i.e.,
$y_D$) with the accumulated value ($T_\mathsf{Acc}$) of $AS$.  During the
authentication phase, $D$ uses $y_D$ again to prove its group membership.
PUF responses and a session-based random are also exchanged to secure the
communication of $y_D$ during authentication.  Additionally, the phase
initialises the pseudonym and dual-state records that protect $D$'s identity
on all subsequent open-channel exchanges.  This phase is assumed to be
executed once, over an end-to-end encrypted channel, for every newly joined
node.  Figure~\ref{fig:enrolment} details the message flow.

\begin{figure}[!t]
  \centering
  \includegraphics[width=0.9\columnwidth]{fig_enrollment}
  \caption{Enrolment Phase via secure channel.}
  \label{fig:enrolment}
\end{figure}

$D$ initiates the phase by sending a registration request along with its
unique identity $ID_D$ to $AS$—the only time $ID_D$ appears in any network
message.  As a response, $AS$ generates a challenge $c_D \in \mathbb{C}_D$
and a session-based random $m_D$, and forwards them securely to $D$.  Both
credentials secure the communication during authentication. Specifically, $c_D$ enables
$D$ to regenerate its PUF response, while $m_D$ introduces unpredictability
into open-channel messages, preventing adversarial prediction.  With $c_D$
as input to its embedded PUF, $D$ obtains $R_D = \mathsf{PUF}_D(c_D)$.
$D$ also generates its unique secret $y_D$ (a random value), and the security
of the entire scheme critically depends on maintaining the confidentiality
of $y_D$.  $D$ picks a reverse challenge $c_{AS\text{-}D}$ and sends
$(y_D, R_D, c_{AS\text{-}D})$ back to $AS$ over the secure channel,
maintaining the confidentiality of $y_D$.  $D$ stores
$\{y_D,\, c_D,\, m_{\mathrm{curr}}{=}m_D\}$.

$AS$ inputs $y_D$ to the hash function $H(\cdot)$, generating the tag
$Y_D^H = H(y_D)$, then performs a bitwise AND with the current accumulated
value, obtaining $T_\mathsf{Acc} = T_\mathsf{Acc} \mathbin{\&} Y_D^H$.  The
accumulated value contains the compact, hashed, unique secrets of all
enrolled nodes, which are authenticated by $AS$.  In this way, $D$
successfully enrols itself with $AS$.  $AS$ then generates
$R_{AS\text{-}D} = \mathsf{PUF}_{AS}(c_{AS\text{-}D})$ and calculates
$\Phi_{AS\text{-}D} = R_{AS\text{-}D} \oplus R_D$, keeping $R_D$ in
XOR-protected form for further secure communications.

To protect $D$'s identity in all subsequent open-channel messages, $AS$
initialises a per-session pseudonym record.  Since $ID_D$ must not appear
on the open channel after enrolment, $D$ will instead be identified by the
\emph{pseudonym} $\mathit{PID} = H(ID_D \| m_{\mathrm{curr}})$ in all future
authentication messages.  Because $H$ is one-way and $m_{\mathrm{curr}}$ is
private, this pseudonym cannot be linked to $ID_D$ or to pseudonyms from
other sessions by a channel observer.  $AS$ pre-computes
$\mathit{PID}_{\mathrm{curr}} = H(ID_D \| m_{\mathrm{curr}})$ and sets
$\mathit{PID}_{\mathrm{old}} = \perp$, $m_{\mathrm{old}} = \perp$.  The
\emph{dual state} ($m_{\mathrm{curr}}$, $m_{\mathrm{old}}$,
$\mathit{PID}_{\mathrm{curr}}$, $\mathit{PID}_{\mathrm{old}}$) is maintained
because the Key Exchange Phase delivers the updated session random $m_{\mathrm{new}}$ over
the open channel, where a single packet loss can cause $D$ to retain
$m_{\mathrm{curr}}$ while $AS$ has already advanced.  By retaining both the
current and the previous pseudonym, $AS$ can recover from such
desynchronisation without requiring $D$ to re-enrol.  $AS$ stores
$\{T_\mathsf{Acc},\, ID_D,\, \Phi_{AS\text{-}D},\, c_{AS\text{-}D},\,
m_{\mathrm{curr}},\, m_{\mathrm{old}},\, \mathit{PID}_{\mathrm{curr}},\,
\mathit{PID}_{\mathrm{old}}\}$.

\subsection{Node Authentication Phase}

The main purpose of this phase is to prove the membership of $D$'s unique
secret $y_D$ to $AS$.  Since authentication is executed over the open
channel, $y_D$ must be communicated in a secure form.  Both the PUF-based
challenge-response, producing the secret $R_D$, and the session-based random
$m_{\mathrm{curr}}$, incorporating unpredictability, are known only to $D$ and
$AS$.  Additionally, the proposed scheme replaces the real identity $ID_D$
with the current pseudonym $\mathit{PID}$, so that a passive observer cannot
link successive sessions to the same device.
Figure~\ref{fig:auth} illustrates the phase.

\begin{figure}[!t]
  \centering
  \includegraphics[width=0.9\columnwidth]{fig_auth.jpg}
  \caption{Authentication Phase via open channel.}
  \label{fig:auth}
\end{figure}

The phase is initiated by $D$.  $D$ regenerates $R_D = \mathsf{PUF}_D(c_D)$
from the stored challenge $c_D$, obtained during enrolment.  Rather than
including its real identity $ID_D$—which would enable cross-session
tracking—$D$ computes its current pseudonym $\mathit{PID} = H(ID_D \| m_{\mathrm{curr}})$.
$D$ then produces the hashed secret $Y_D^H = H(y_D)$ in XOR-masked form as
$Y_{AS\text{-}D}^H = H(y_D) \oplus H(R_D \| m_{\mathrm{curr}} \| \mathit{PID} \| ts_1)$.
The XOR mask is the hash of $R_D$ (regenerated from $c_D$), the
session-based random $m_{\mathrm{curr}}$, the pseudonym $\mathit{PID}$, and
the current timestamp $ts_1$.  The timestamp provides replay protection to
the XOR-masked $Y_D^H$.  The complete message $(\mathit{PID},\,
Y_{AS\text{-}D}^H,\, ts_1)$ is sent to $AS$.  An attacker cannot obtain
$H(y_D)$ from $Y_{AS\text{-}D}^H$ without knowing $R_D$ (which can only be
generated by $D$'s PUF) and $m_{\mathrm{curr}}$ (known only to $D$ and $AS$).

After receiving the message, $AS$ performs a \emph{dual-state lookup}~\cite{zhong2025survey}, checking
whether the received $\mathit{PID}$ matches $\mathit{PID}_{\mathrm{curr}}$
or $\mathit{PID}_{\mathrm{old}}$ in its records.  A match on
$\mathit{PID}_{\mathrm{old}}$ indicates that $D$ missed the Key Exchange Phase delivery
in the previous session and still holds the prior $m_{\mathrm{curr}}$. In this
case, $AS$ uses the stored $m_{\mathrm{old}}$ for this verification, recovering
from the desynchronisation without requiring re-enrolment.  $AS$ verifies
the freshness of $ts_1$.  $R_D'$ is calculated as $\Phi_{AS\text{-}D} \oplus
\mathsf{PUF}_{AS}(c_{AS\text{-}D})$, where both $\Phi_{AS\text{-}D}$ and
$c_{AS\text{-}D}$ are retrieved from $D$'s enrolment record.  $AS$
regenerates the same XOR mask and unmasks $Y_D^{H'}$, then performs the
bitwise AND check: $T_\mathsf{Acc}^\mathsf{new} = T_\mathsf{Acc} \mathbin{\&} Y_D^{H'} \stackrel{?}{=} T_\mathsf{Acc}$.
From the properties of the hash-based one-way accumulator, if an element is
combined into an accumulated value using bitwise AND and the same operation is
again performed between them, it gives back the same accumulated value.  If
$T_\mathsf{Acc}^\mathsf{new} = T_\mathsf{Acc}$, it proves that the unique
secret of $D$ is already a legitimate member of $AS$, completing
authentication.  On success, $AS$ proceeds to the Key Exchange Phase.

\subsection{Node Session Key Exchange Phase}

This phase is executed by $AS$ to inform $GW$ about the successful
authentication of the newly joined node $D$ and to mediate the session key
between $GW$ and $D$.  The proposed scheme additionally derives a fresh pseudonym $\mathit{PID}_{\mathrm{new}}$ for
$D$ at the end of each session and updates the dual-state records at both
$AS$ and $GW$, ensuring that successive sessions are unlinkable.  The two
deliveries to $D$ and to $GW$ are sent simultaneously.
Figure~\ref{fig:keyex} shows the phase.

\begin{figure}[!t]
  \centering
  \includegraphics[width=0.9\columnwidth]{fig_keyex}
  \caption{Session Key Exchange Phase via open channel.}
  \label{fig:keyex}
\end{figure}

$AS$ uses $\mathsf{auth\_token}_D$ to deliver the session key $K_{GW\text{-}D}$
to $GW$.  $K_{GW\text{-}D}$ is computed as $H(R_D' \| m_{\mathrm{new}})$,
where $R_D'$ is the PUF response of $D$ regenerated by $AS$ during the
authentication phase and thus known only to $AS$ and $D$, and
$m_{\mathrm{new}} = H(n_1)$ is derived from a fresh random $n_1$ generated by
$AS$.  The computed $m_{\mathrm{new}}$ also serves as the updated session-based
random for the next authentication session between $AS$ and $D$,
incorporating an update to the authentication secret.  In addition, $AS$
derives the session key $K_{GW\text{-}D} = H(R_D' \| m_{\mathrm{new}})$ and
the next pseudonym $\mathit{PID}_{\mathrm{new}} = H(ID_D \| m_{\mathrm{new}})$.
Because $\mathit{PID}_{\mathrm{new}}$ is based on the fresh $m_{\mathrm{new}}$,
it is computationally unlinkable to the $\mathit{PID}$ used in the current
session.  The session key, $ID_{AS}$, $ID_D$, the next pseudonym
$\mathit{PID}_{\mathrm{new}}$, and the authentication completion timestamp
$ts_{\mathrm{auth}}$ are encrypted under $K_{GW\text{-}AS}$ to form
$\mathsf{auth\_token}_D$.  The token carries $\mathit{PID}_{\mathrm{new}}$ so
that $GW$ can update its pseudonym mapping for $D$'s future communications.

$AS$ also informs $D$ of $m_{\mathrm{new}}$ so that $D$ can independently
derive $K_{GW\text{-}D}$ and its next pseudonym.  $AS$ masks $m_{\mathrm{new}}$
to protect it from adversaries—if an adversary captures $m_{\mathrm{new}}$, the
probability of predicting future session key material increases—so it is
masked as $m_H = m_{\mathrm{new}} \oplus H(Y_D^{H'} \| m_{\mathrm{curr}} \| R_D' \| ID_{AS} \| \mathit{PID}_{\mathrm{curr}} \| ts_2)$.
Here, $Y_D^{H'}$ is the hashed group secret known to $AS$ from the
authentication phase, $m_{\mathrm{curr}}$ and $R_D'$ are likewise known from
authentication, and $ts_2$ prevents replay of $m_H$ in a future session.
$AS$ simultaneously sends $(m_H,\, ts_2)$ to $D$ and $(ID_{AS},\,
\mathsf{auth\_token}_D)$ to $GW$.  It then saves the current nonce and
pseudonym as the previous values and updates its records with
$m_{\mathrm{new}}$ as the current nonce and $\mathit{PID}_{\mathrm{new}}$
as the current pseudonym.

Obtaining $(m_H,\, ts_2)$, $D$ first verifies the freshness of $ts_2$
(freshly generated by $AS$ and included in the delivery), then recovers
$m_{\mathrm{new}}$ by regenerating the hash mask using its locally held
$Y_D^H$, $R_D$, $m_{\mathrm{curr}}$, $ID_{AS}$, $\mathit{PID}$, and the
received $ts_2$, and XOR-ing with $m_H$.  $D$ establishes the session key
$K_{GW\text{-}D} = H(R_D \| m_{\mathrm{new}})$, updates its pseudonym
$\mathit{PID}_{\mathrm{curr}} = H(ID_D \| m_{\mathrm{new}})$, updates the
stored session nonce to $m_{\mathrm{new}}$, and stores $K_{GW\text{-}D}$.

$GW$ decrypts the token with $K_{GW\text{-}AS}$, verifies $ts_{\mathrm{auth}}$,
and extracts $K_{GW\text{-}D}$ and $\mathit{PID}_{\mathrm{new}}$.  $GW$ saves
$\mathit{PID}_{\mathrm{curr}}$ as the previous pseudonym, sets
$\mathit{PID}_{\mathrm{new}}$ as the current one, and stores the session
key against both.
Retaining $\mathit{PID}_{\mathrm{old}}$ ensures that if the Key Exchange Phase delivery
to $D$ was lost, $GW$ can still service $D$'s data frames under the prior
session key pending re-authentication.

\subsection{Desynchronization Recovery}

If the Key Exchange Phase delivery $(m_H,\,ts_2)$ is lost, $D$ retains its stale
$m_{\mathrm{curr}}$, so the pseudonym it computes in the next session equals
$\mathit{PID}_{\mathrm{old}}$ from $AS$'s perspective.  Algorithm~\ref{alg:desync}
shows how the scheme adapts Phases~2 and~3 to recover without re-enrolment.
$D$'s behaviour is \emph{identical} in both paths—it always uses its own
$m_{\mathrm{curr}}$ and $\mathit{PID}$.  Only $AS$ adapts: the dual-state
lookup detects which nonce $D$ is on, then constructs $m_H$ with the matching
$(m_{\mathrm{old}},\,\mathit{PID}_{\mathrm{old}})$ so that $D$'s symmetric
decryption succeeds and the states converge.
The dual-state mechanism tolerates a single consecutive Key Exchange Phase loss, but
two or more consecutive losses exhaust both slots, requiring re-enrolment---a
deliberate trade-off between storage overhead and recovery depth.

\begin{algorithm}[t]
\caption{Desynchronization Recovery}
\label{alg:desync}
\begin{algorithmic}[1]
\Statex \textbf{Precondition:} $D$ holds stale $m_{\mathrm{curr}}\!=\!m_{\mathrm{old}}$;\;
       $AS$ holds $(m_{\mathrm{curr}},\,m_{\mathrm{old}},\,
       \mathit{PID}_{\mathrm{curr}},\,\mathit{PID}_{\mathrm{old}})$

\Statex
\Statex \textit{Authentication Phase}
\State $\mathit{PID} \gets H(\mathit{ID}_D \| m_{\mathrm{curr}})$
       \hfill\Comment{stale $m_{\mathrm{curr}}$ $\Rightarrow$ $\mathit{PID} = \mathit{PID}_{\mathrm{old}}$}
\State $D \to AS$: $(\mathit{PID},\; Y_{AS\text{-}D}^H,\; ts_1)$
\If{$\mathit{PID} = \mathit{PID}_{\mathrm{curr}}$}
    \State $AS$ verifies using $m_{\mathrm{curr}}$ \hfill\Comment{normal path}
\ElsIf{$\mathit{PID} = \mathit{PID}_{\mathrm{old}}$}
    \State $AS$ verifies using $m_{\mathrm{old}}$ \hfill\Comment{desync detected}
\Else
    \State $AS$ rejects $D$
\EndIf

\Statex
\Statex \textit{Key Exchange Phase (desync path)}
\State $m_{\mathrm{new}} \gets \mathsf{rand}()$
\State $m_H \gets m_{\mathrm{new}} \oplus
       H\!\left(Y_D^{H'} \| m_{\mathrm{old}} \| R_D' \|
       \mathit{ID}_{AS} \| \mathit{PID}_{\mathrm{old}} \| ts_2\right)$
       \hfill\Comment{$AS$ uses $m_{\mathrm{old}}$, $\mathit{PID}_{\mathrm{old}}$}
\State $AS \to D$ (via $GW$): $(m_H,\; ts_2)$
\State $m_{\mathrm{new}} \gets m_H \oplus
       H\!\left(Y_D^{H} \| m_{\mathrm{curr}} \| R_D \|
       \mathit{ID}_{AS} \| \mathit{PID} \| ts_2\right)$
       \hfill\Comment{$D$ uses its own $m_{\mathrm{curr}}\!=\!m_{\mathrm{old}}$ — recovers correctly}
\State $m_{\mathrm{old}} \gets m_{\mathrm{curr}}$;\quad
       $m_{\mathrm{curr}} \gets m_{\mathrm{new}}$;\quad
       $\mathit{PID} \gets H(\mathit{ID}_D \| m_{\mathrm{curr}})$
\Statex \textbf{Result:} $D$ and $AS$ re-synchronised; $\mathit{PID}$ rotated
\end{algorithmic}
\end{algorithm}

% -------------------------------------------------------
\FloatBarrier
\section{Security Analysis}

\subsection{Formal Security Analysis}

The proposed scheme is translated into ProVerif~\cite{blanchet2018proverif}
queries under the Dolev--Yao adversary model~\cite{dolev1983security}, where
the adversary has full control of the public channel and can intercept, replay,
modify, and inject messages.  Ten queries are verified and all are satisfied.
Figure~\ref{fig:proverif} shows the verification summary.

\begin{figure*}[!t]
  \centering
  \includegraphics[width=\textwidth]{fig_proverif}
  \caption{ProVerif Output.}
  \label{fig:proverif}
\end{figure*}

\subsubsection{Protocol Correspondence}

Five correspondence queries verify that the protocol steps happen in the correct
order and that no adversary can skip or forge a step.

The first query checks enrollment integrity: a device can only be registered at
the authentication server after it has explicitly initiated the enrollment
process.  This rules out any possibility of a device being enrolled without
its own participation.

The second query checks the device side of the authentication round: the device
receives the server's acknowledgement and the fresh timestamp $ts_2$ only after
it first sent a valid authentication request.  This prevents an adversary from
replaying a stale acknowledgement to trick the device into advancing its state.

The third query checks the server side of the same round: the authentication
server completes its processing only after it received a structurally correct
request from the device.  A forged or replayed message cannot push the server
into a completed state.

The fourth query binds the two rounds together: the PUF response $R_D$ and
the session seed $m_D$ that the device commits in Round~1 are proven to be
exactly the same values that appear in Round~2.  This rules out an active
adversary who might try to substitute credentials between rounds.

The fifth query covers key exchange: the device recovers the new session seed
$m_{\mathrm{new}}$ only after it legitimately sent the key-exchange request.
No key material can be delivered to the device out of sequence.

\subsubsection{Secrecy and Forward Secrecy}

Two queries verify the secrecy of the session key
$K_{GW\text{-}D} = H(R_D \| m_{\mathrm{new}})$ from the perspective of the
device and the gateway independently.  Even though the adversary observes all
traffic on the public channel, it cannot reconstruct $K_{GW\text{-}D}$
because the PUF response $R_D$ is never transmitted and
$m_{\mathrm{new}}$ is always delivered encrypted.

A separate query confirms forward secrecy: the rotated session seed
$m_{\mathrm{new}}$ itself remains secret.  This means that even if an
adversary somehow recovers a current session key, it gains nothing about
future sessions, because each $m_{\mathrm{new}}$ is derived from an
independently generated fresh nonce.

\subsubsection{Anonymity Preservation}

The most significant secrecy result concerns the real device identity $ID_D$.
After enrollment, $ID_D$ is never sent over the open channel again.  Every
subsequent message instead carries the pseudonym
$\mathit{PID} = H(ID_D \| m_{\mathrm{curr}})$.
ProVerif formally confirms that no attacker, even one who intercepts and manipulates every single
message on the network, can recover $ID_D$ from the observed traffic.

This result supports device anonymity.  However, anonymity is only meaningful
if the pseudonym used in one session cannot be linked to the pseudonym used in
the next.  The forward-secrecy query addresses this: because
$m_{\mathrm{new}}$ is proven secret, consecutive pseudonyms
$H(ID_D \| m_{\mathrm{curr}})$ and $H(ID_D \| m_{\mathrm{new}})$ are
computationally indistinguishable to any observer, making it difficult for an
adversary who collects pseudonyms from multiple sessions to determine whether
they belong to the same physical device.

\subsubsection{Offline Guessing Resistance}

Finally, a weak-secrecy query confirms that the session key $K_{GW\text{-}D}$
cannot be guessed from the public transcript even by an adversary with
unlimited offline dictionary power.  This closes the gap against brute-force
and dictionary attacks that target the key directly from captured packets.

\subsection{Informal Analysis}

\subsubsection{Replay Attack Resistance}
The Enrolment Phase is executed over a secure channel and immune to replay.
The Authentication and Key Exchange Phases embed timestamps
($ts_1$, $ts_2$, $ts_{\mathrm{auth}}$) and the session-based random
$m_{\mathrm{curr}}$ into the hash masks.  A replayed message with a stale
timestamp or stale $m_{\mathrm{curr}}$ will fail the freshness check at
$AS$, $D$, or $GW$ respectively.

\subsubsection{MITM Attack Resistance}
The Enrolment Phase is conducted over an encrypted channel, which precludes
MITM.  In the Authentication and Key Exchange Phases, the PUF response $R_D$ (device-specific, unreproducible) and
session random $m_{\mathrm{curr}}$ (private to $D$ and $AS$) are embedded in
every XOR mask.  Without $R_D$ and the current $m_{\mathrm{curr}}$, no
adversary can forge a valid masked tag $Y_{AS\text{-}D}^H$ or unmask $m_H$.

\subsubsection{Device Anonymity}
$ID_D$ is transmitted only once, over the secure enrolment channel, and is
never exposed on the open network again.  All open-channel messages use
$\mathit{PID} = H(ID_D \| m_{\mathrm{curr}})$, which is refreshed to
$H(ID_D \| m_{\mathrm{new}})$ after every key exchange.  Because $m_{\mathrm{new}}$
is derived from a freshly generated nonce $n_1$, successive pseudonyms are
computationally unlinkable.

\subsubsection{Desynchronization Recovery}
Packet loss can occur at three distinct points in Phases~2 and~3, each
handled without requiring re-enrolment.

If the acknowledgement and fresh timestamp $ts_2$ sent by $AS$ to $D$ at the
end of Round~1 are lost in transit, $D$ does not advance its internal state
and retains $m_{\mathrm{curr}}$ and the current $\mathit{PID}$ unchanged.
Since $AS$ has not yet performed any state update at this stage, both
entities remain consistent.  Upon timeout, $D$ retransmits the Round~1
request with the same $\mathit{PID}$ and $ts_1$, and the exchange proceeds
normally.

The primary desynchronization scenario arises when the Key Exchange Phase response from
$AS$ to $D$ — carrying the masked session seed $m_H$ — is lost.  In this
case, $AS$ has already advanced its state to $m_{\mathrm{new}}$ and rotated
the pseudonym to $\mathit{PID}_{\mathrm{new}}$, while $D$ retains the
previous $m_{\mathrm{curr}}$.  On the subsequent authentication attempt, $D$
constructs its pseudonym from the stale $m_{\mathrm{curr}}$, producing
$\mathit{PID}_{\mathrm{old}}$.  $AS$ recognises this via the dual-state
lookup — storing both $(m_{\mathrm{curr}}, \mathit{PID}_{\mathrm{curr}})$ and
$(m_{\mathrm{old}}, \mathit{PID}_{\mathrm{old}})$ — and resumes from the
consistent prior state, restoring synchronisation without any re-enrolment
overhead.

If the token forwarded from $AS$ to $GW$ is lost, $GW$ does not receive
$\mathit{PID}_{\mathrm{new}}$ or the associated session key $K_{GW\text{-}D}$.
Consequently, $D$'s subsequent data packet, carrying $\mathit{PID}_{\mathrm{new}}$,
is not recognised by $GW$.  $D$ re-initiates from the Authentication Phase; the dual-state
storage at $AS$ ensures that this re-authentication succeeds and a fresh
token is forwarded to $GW$, restoring normal operation.

\subsubsection{Forward Secrecy}
$K_{GW\text{-}D} = H(R_D \| m_{\mathrm{new}})$.  Compromise of a session key
does not reveal prFig. 4.   ProVerif Output.evious session keys because $m_{\mathrm{new}}$ is
independently generated from a fresh nonce and $R_D$ is PUF-derived.

% -------------------------------------------------------
\section{Performance Analysis}

\subsection{Theoretical Performance Analysis}
\subsubsection{Computational Cost}

The computational cost is a crucial deterministic factor ensuring the security and
efficiency of an authentication scheme. Pycrypto, a Python cryptographic library, is used
in a Windows system to measure the computational cost of core cryptographic functions,
such as SHA-256 ($T_{\mathsf{hash}}$), AES Encryption/Decryption ($T_{\mathsf{aes}}$),
and Random generation ($T_{\mathsf{rand}}$). The PUF operation cost ($T_{\mathsf{puf}}$)
is taken from prior benchmarks on the same platform~\cite{zheng2022puf}.

To execute the node enrolment phase (Figure~\ref{fig:enrolment}):
\begin{itemize}
  \item $AS$ computes two random generations ($T_{\mathsf{rand}}$), one PUF operation
        ($T_{\mathsf{puf}}$), and two hash operations ($T_{\mathsf{hash}}$) for pseudonym
        initialisation.
  \item On the other hand, $D$ executes two random generations ($T_{\mathsf{rand}}$) and
        one PUF operation ($T_{\mathsf{puf}}$).
\end{itemize}

Hence, the total computational cost incurred per node enrolment amounts to
$2T_{\mathsf{puf}} + 2T_{\mathsf{hash}} + 4T_{\mathsf{rand}}$.
As shown in Table~\ref{tab:comp_enroll}, the proposed scheme achieves an enrolment cost of
0.18\,ms. It competes well with Zhou~\cite{zhou2024auth} (0.13\,ms) while
additionally providing PUF-based secret binding and pseudonym initialisation, both of
which are absent from those schemes.

\begin{table}[!t]
  \centering
  \caption{Computational Cost Comparison: Enrolment Phase.}
  \label{tab:comp_enroll}
  \scriptsizeFig. 4.   ProVerif Output.
  \setlength{\tabcolsep}{3pt}
  \begin{tabular}{lll}
    \toprule
    Scheme & Operations & Cost (ms) \\
    \midrule
    LAAKA~\cite{ali2024laaka}      & $3T_\mathsf{hash}$ & 0.11 \\
    Zhou~\cite{zhou2024auth}      & $T_\mathsf{puf}+2T_\mathsf{hash}+3T_\mathsf{rand}$ & 0.13 \\
    Zheng et al.~\cite{zheng2022puf}      & $2T_\mathsf{puf}+8T_\mathsf{hash}+4T_\mathsf{aes}+5T_\mathsf{rand}$ & 0.69 \\
    He~et~al.~\cite{he2023lightweight}  & $4T_\mathsf{hash}+T_\mathsf{aes}+T_\mathsf{rand}$ & 0.23 \\
    \textbf{Proposed}        & $2T_\mathsf{puf}+2T_\mathsf{hash}+4T_\mathsf{rand}$ & \textbf{0.18} \\
    \bottomrule
  \end{tabular}\\[3pt]
  \scriptsize
  Benchmark (avg.\ 100 iterations): $T_\mathsf{rand}$=0.0018\,ms, $T_\mathsf{puf}$=0.0522\,ms,
  $T_\mathsf{hash}$=0.0353\,ms, $T_\mathsf{aes}$=0.0867\,ms.
\end{table}

Due to the execution of the authentication and key exchange phase, analysis of its
computational efficiency is essential. In the proposed scheme, $D$ and $AS$
(Figures~\ref{fig:auth} and~\ref{fig:keyex}) both execute one PUF operation
($T_{\mathsf{puf}}$).
\begin{itemize}
  \item $D$ additionally performs six hash operations ($T_{\mathsf{hash}}$): pseudonym
        computation, group secret hash, and XOR mask in the Authentication Phase, followed by mask
        recovery, session key derivation, and pseudonym rotation in the Key Exchange Phase.
  \item $AS$ performs five hash operations ($T_{\mathsf{hash}}$), along with one AES
        encryption ($T_{\mathsf{aes}}$) for the gateway token and one random generation
        ($T_{\mathsf{rand}}$) for $m_{\mathrm{new}}$.
  \item $GW$ runs one AES decryption ($T_{\mathsf{aes}}$) to recover the session key and
        updated pseudonym from the forwarded token.
\end{itemize}

This leads to the total computational cost of the authentication and key exchange phase
amounting to $2T_{\mathsf{puf}} + 11T_{\mathsf{hash}} + 2T_{\mathsf{aes}} + 1T_{\mathsf{rand}}$.
In Table~\ref{tab:comp_auth}, the proposed scheme is said to have a computational cost of
0.67\,ms. Its computational efficiency surpasses that of
He~et~al.~\cite{he2023lightweight} (1.74\,ms), whose ECC scalar multiplication
operations incur significantly higher cost on constrained platforms. Due to the incorporation of a PUF
challenge--response pair and an AES-encrypted gateway token — enabling PUF-based group
membership verification and decoupled AS selection — the scheme incurs a modest overhead
over LAAKA~\cite{ali2024laaka} (0.56\,ms) and Zhou~\cite{zhou2024auth} (0.53\,ms), which
neither employ PUF-based group membership, a decoupled distributed AS, nor dual-state
desynchronisation recovery.

\begin{table}[!t]
  \centering
  \caption{Computational Cost Comparison: Authentication \& Key Exchange Phase.}
  \label{tab:comp_auth}
  \scriptsize
  \setlength{\tabcolsep}{3pt}
  \begin{tabular}{lll}
    \toprule
    Scheme & Operations & Cost (ms) \\
    \midrule
    LAAKA~\cite{ali2024laaka}      & $16T_\mathsf{hash}$ & 0.56 \\
    Zhou~\cite{zhou2024auth}      & $15T_\mathsf{hash}$ & 0.53 \\
    Zheng et al.~\cite{zheng2022puf}      & $2T_\mathsf{puf}+6T_\mathsf{hash}+4T_\mathsf{aes}+5T_\mathsf{rand}$ & 0.67 \\
    He~et~al.~\cite{he2023lightweight}  & $11T_\mathsf{hash}+5T_\mathsf{aes}+2T_\mathsf{ecc}$ & 1.74 \\
    \textbf{Proposed}        & $2T_\mathsf{puf}+11T_\mathsf{hash}+2T_\mathsf{aes}+T_\mathsf{rand}$ & \textbf{0.67} \\
    \bottomrule
  \end{tabular}\\[3pt]
  \scriptsize
  $T_\mathsf{ecc}$=0.461\,ms (scaled from He~et~al.~\cite{he2023lightweight} hardware ratios).
\end{table}

\subsubsection{Communication Cost}

Communication cost is a core feature for analysing the performance of an authentication
scheme, based on the number of message exchanges and their sizes (in bits). By assuming
that communications are executed with uniform bit lengths, the communication cost of both
enrolment and the authentication and key exchange phase is evaluated and compared with
other state-of-the-art schemes.

It is assumed that identities and timestamps are 32\,bits, along with nonces, hashes,
challenges, responses, and ciphertexts as 256\,bits, ensuring a fair comparison.

During the node enrolment phase, $D$ enrols itself with $AS$, executing three message
exchanges. The first message communicates the unique identity of $D$, i.e., $ID_D$
(32\,bits), to $AS$ over the secure channel. In the second message, $AS$ transmits a
challenge $c_D$ (256\,bits) and a session-based random $m_D$ (256\,bits) to $D$. The
third communication takes place from $D$ to $AS$ in the form of a unique secret $y_D$
(256\,bits), a PUF response $R_D$ (256\,bits), and a reverse challenge
$c_{AS\text{-}D}$ (256\,bits). Thus, the total communication cost per enrolment is
1312\,bits as presented in Table~\ref{tab:comm_enroll}. It outperforms schemes such
as LAAKA~\cite{ali2024laaka} (2592\,bits) and Zheng et al.~\cite{zheng2022puf} (1600\,bits), while remaining
competitive with Zhou~\cite{zhou2024auth} (1344\,bits) and He~et~al.~\cite{he2023lightweight} (1280\,bits).
Unlike those schemes, the proposed scheme binds device enrolment to a PUF-derived secret and initialises
the pseudonym at registration, adding security guarantees at a modest communication cost.

\begin{table}[H]
  \centering
  \caption{Communication Cost Comparison: Enrolment Phase.}
  \label{tab:comm_enroll}
  \scriptsize
  \setlength{\tabcolsep}{3pt}
  \begin{tabular}{lcc}
    \toprule
    Scheme & \# Messages & Cost (bits) \\
    \midrule
    LAAKA~\cite{ali2024laaka}      & 3 & 2592 \\
    Zhou~\cite{zhou2024auth}      & 5 & 1344 \\
    Zheng et al.~\cite{zheng2022puf}      & 3 & 1600 \\
    He~et~al.~\cite{he2023lightweight}  & 2 & 1280 \\
    \textbf{Proposed}        & 3 & \textbf{1312} \\
    \bottomrule
  \end{tabular}\\[3pt]
  \scriptsize
  IDs/timestamps = 32\,b; hashes/nonces/challenges/ciphertexts = 256\,b.
\end{table}
\begin{table}[H]
  \centering
  \caption{Communication Cost Comparison: Auth.\ \& Key Exchange Phase.}
  \label{tab:comm_auth}
  \scriptsize
  \setlength{\tabcolsep}{3pt}
  \begin{tabular}{lcc}
    \toprule
    Scheme & \# Messages & Cost (bits) \\
    \midrule
    LAAKA~\cite{ali2024laaka}      & 3 & 2400 \\
    Zhou~\cite{zhou2024auth}      & 4 & 2464 \\
    Zheng et al.~\cite{zheng2022puf}      & 3 & 1600 \\
    He~et~al.~\cite{he2023lightweight}  & 3 & 1120 \\
    \textbf{Proposed}        & 3 & \textbf{1376} \\
    \bottomrule
  \end{tabular}
\end{table}
The authentication and key exchange phase is executed using three message communications.
The first message takes place from $D$ to $AS$, where $D$ communicates its current
pseudonym $\mathit{PID}$ (256\,bits), a masked membership tag $Y_{AS\text{-}D}^H$
(256\,bits), and a timestamp $ts_1$ (32\,bits). In the second message, $AS$ transmits the
masked session seed $m_H$ (256\,bits) and a timestamp $ts_2$ (32\,bits) to $D$. The
third communication takes place from $AS$ to $GW$, consisting of the updated pseudonym
$\mathit{PID}_{\mathrm{new}}$ (256\,bits), the identity $ID_{AS}$ (32\,bits), and the
encrypted gateway token $\mathsf{auth\_token}_D$ (256\,bits). Thus, the total
communication cost becomes 1376\,bits, as presented in Table~\ref{tab:comm_auth}.
He~et~al.~\cite{he2023lightweight} achieves a lower figure (1120\,bits) for a simpler
two-party IoT-to-server protocol, whereas the proposed scheme supports multihop mesh
networks with a decoupled distributed AS, desynchronisation recovery, and PUF-based
membership verification — features that account for the additional overhead.
Das~\cite{das2026comsnets} requires 928\,bits, with the 448-bit increase in the proposed
scheme directly reflecting the widening of the device identifier from a 32-bit $ID_D$ to
a 256-bit pseudonym $\mathit{PID}$ — the direct cost of adding device anonymity to all
open-channel communications.




\subsection{Simulation-Based Performance Analysis}
\label{sec:sim}

A simulation-based performance evaluation has been conducted to detect the efficiency of
the proposed autheFig. 4.   ProVerif Output.ntication scheme using the COOJA simulator within the Contiki-NG
framework. Three complementary studies are reported. The first is an overall per-device
comparison of the proposed scheme against LAAKA~\cite{ali2024laaka} and
Zhou~\cite{zhou2024auth} in a fixed 100-node network comprising 1 gateway, 79 AS nodes
(2 active), and 20 newly joined devices. The second is an authentication server (AS)
pool-size study that varies the number of active AS nodes from 2 to 10. The third is a
network scalability study that grows the total node count from 50 to 100. All results are averaged over five independent
random seeds. Table~\ref{tab:simparams} elaborates the common simulation parameters.

\begin{table}[H]
  \centering
  \caption{Simulation Parameters.}
  \label{tab:simparams}
  \scriptsize
  \setlength{\tabcolsep}{4pt}
  \begin{tabular}{ll}
    \toprule
    Parameter & Value \\
    \midrule
    Simulator          & COOJA (Contiki-NG) \\
    Mote type          & TelosB (emulated) \\
    Network topology (study~i) & 100 nodes: 1 GW, 79 AS (2 active), 20 devices \\
    Phases measured    & Enrolment, Authentication, Key Exchange \\
    Schemes compared   & Proposed scheme, LAAKA~\cite{ali2024laaka}, Zhou~\cite{zhou2024auth} \\
    Active AS (study~ii) & 2, 5, 10 \\
    Network size (study~iii) & $N \in \{50, 80, 100\}$ nodes \\
    Random seeds       & 5 per configuration \\
    \bottomrule
  \end{tabular}
\end{table}

\subsubsection{Performance Analysis for Overall Per-Device Cost}

Initially, the authentication schemes are evaluated by considering the overhead of one
complete authentication cycle — enrolment, authentication, and key exchange — per device.
For 20 newly joined nodes in a fixed 100-node network (1 gateway, 79 AS nodes with 2 active,
and 20 newly joined devices), the proposed scheme demonstrated a
consistent advantage over LAAKA~\cite{ali2024laaka} and Zhou~\cite{zhou2024auth} by
achieving a reduction of 42.8\,\% and 34.6\,\% with respect to total per-device energy
consumption (Figure~\ref{fig:sim_total}a), respectively. Similarly, a reduction of
42.7\,\% and 35.3\,\% is demonstrated with respect to total per-device CPU time
(Figure~\ref{fig:sim_total}b) over LAAKA and Zhou, respectively. The dominant savings
arise in the authentication and key exchange phase, where the proposed scheme consumes
28.87\,mJ against 77.97\,mJ for LAAKA and 57.06\,mJ for Zhou. Due to the incorporation
of fewer message exchanges and computationally lightweight hash-XOR operations, the
proposed scheme has minimised radio activity and computational load compared to the
message-intensive and cryptographically heavier LAAKA and Zhou schemes.

\begin{figure}[!t]
  \centering
  \includegraphics[width=\columnwidth]{fig_sim_total}
  \caption{Per-device mean total cost — all phases:
           (a)~energy and (b)~CPU time.}
  \label{fig:sim_total}
\end{figure}

\subsubsection{Performance Analysis for Authentication Server Pool Size}

In the decoupled design, the gateway selects any capable node as the authentication
server, so the size of the AS pool directly influences load distribution and overall
network performance. The performance of the proposed scheme, LAAKA, and Zhou is evaluated
by varying the number of active AS nodes from 2 to 10. Considering the total energy
consumption (Figure~\ref{fig:sim_as}a), a reduction of approximately 42\,\% is
demonstrated by the proposed scheme over LAAKA at every AS count. Both schemes improve
from AS\,=\,2 to AS\,=\,10, with the proposed scheme dropping from 1065\,mJ to
936\,mJ and LAAKA from 1827\,mJ to 1631\,mJ, suggesting that a pool of 10 AS nodes is near-optimal
for a 20-device network. Similarly, as shown in Figure~\ref{fig:sim_as}(b), a reduction
of approximately 42\,\% in total CPU time is demonstrated over LAAKA across all AS
counts. Zhou's total cost remains flat at approximately 1503\,mJ in energy and 24.4\,s in
CPU time regardless of AS pool size, confirming that its gateway-coupled design is
insensitive to AS selection diversity. These consistent reductions demonstrate the energy
and computational efficiency of the decoupled AS design in the proposed scheme.

\begin{figure}[!htb]
  \centering
  \includegraphics[width=0.78\columnwidth]{fig_sim_as_energy}\\[2pt]
  \small (a)\\[5pt]
  \includegraphics[width=0.78\columnwidth]{fig_sim_as_cpu}\\[2pt]
  \small (b)
  \caption{Total authentication cost vs.\ number of active AS nodes (20 devices):
           (a)~energy and (b)~CPU time.}
  \label{fig:sim_as}
\end{figure}

\subsubsection{Performance Analysis for Network Scalability}

The performance of authentication schemes is also analysed by considering their
scalability with respect to total network size $N$, varied from 50 to 100 nodes.
Considering the total energy consumption (Figure~\ref{fig:sim_network}a), a reduction of
42.2\,\% and 50.6\,\% is demonstrated by the proposed scheme at $N = 100$ over LAAKA
(898\,mJ) and Zhou (1051\,mJ), respectively. Similarly, as shown in
Figure~\ref{fig:sim_network}(b), a reduction of 42.5\,\% and 50.6\,\% is demonstrated
with respect to total CPU time over LAAKA (14.6\,s) and Zhou (17.0\,s), respectively.
LAAKA and Zhou exhibit steeper energy and CPU growth with increasing $N$ because their
heavier per-device cryptographic load — AES-intensive authentication in LAAKA and a
combined heavy Auth+KeyEx round in Zhou — accumulates faster as the device count grows.
Due to the incorporation of lightweight XOR-hash operations and a decoupled AS design, the
proposed scheme keeps per-device cost low and achieves the lowest total cost among all
compared schemes across all measured network sizes. These consistent reductions across all network
sizes demonstrate the scalability and computational efficiency of the proposed scheme in
comparison to the benchmark schemes.

\begin{figure}[!htb]
  \centering
  \includegraphics[width=0.78\columnwidth]{fig_sim_net_energy}\\[2pt]
  \small (a)\\[5pt]
  \includegraphics[width=0.78\columnwidth]{fig_sim_net_cpu}\\[2pt]
  \small (b)
  \caption{Network scalability: (a)~total energy and
           (b)~total CPU time vs.\ network size~$N$.}
  \label{fig:sim_network}
\end{figure}

% -------------------------------------------------------
% \FloatBarrier
\section{Conclusion}

This paper presented a lightweight, anonymous, and decoupled distributed
authentication scheme for multihop IoT networks.  The assigned
authentication server—which may or may not be the parent of the newly joined
device—performs group-membership verification (Authentication Phase) and session key
establishment (Key Exchange Phase) in two clean rounds.  Device anonymity is preserved
through per-session pseudonym rotation, preventing adversarial tracking across
sessions.  A dual-state mechanism at both the authentication server and the
gateway provides transparent desynchronization recovery without re-enrolment.
PUF-based authentication eliminates long-term secret storage, and
hash-based accumulator operations with XOR masking keep the computational
overhead low.

Formal verification in ProVerif confirms authentication correspondence,
session key secrecy, resistance to offline guessing attacks, and
observational-equivalence-based anonymity.  COOJA/Contiki-NG simulations on emulated TelosB nodes confirm practical
suitability. The proposed scheme achieves a per-device three-phase energy of
52.21\,mJ, reducing total energy by 42.8\,\% over LAAKA~\cite{ali2024laaka}
and 34.6\,\% over Zhou~\cite{zhou2024auth}, with consistent gains across
varying network sizes (up to $N=100$ nodes) and AS pool sizes.

% -------------------------------------------------------
\balance
\bibliographystyle{IEEEtran}
\bibliography{ref}

\iffalse
\begin{thebibliography}{25}

\bibitem{das2026comsnets}
S.~Das and M.~Khatua,
``Lightweight and Decoupled Distributed Authentication for Multihop IoT
Networks,''
in \emph{Proc.\ IEEE COMSNETS}, 2026.

\bibitem{ali2024laaka}
H.~Ali and I.~Ahmed,
``LAAKA: A Lightweight Anonymous Authentication and Key Agreement Scheme for
IoT-Fog Environment,''
\emph{Computers \& Security}, vol.~140, p.~103770, 2024.

\bibitem{zhong2025survey}
J.~Zhong, S.~He, Z.~Liu, and L.~Xiong,
``Lightweight Anonymous Authentication for IoT: A Taxonomy and Survey of
Security Frameworks,''
\emph{Sensors}, vol.~25, no.~17, p.~5594, 2025.

\bibitem{zhou2021iot}
I.~Zhou, I.~Makhdoom, N.~Shariati, M.~A.~Raza, R.~Keshavarz, J.~Lipman,
M.~Abolhasan, and A.~Jamalipour,
``Internet of Things 2.0: Concepts, Applications, and Future Directions,''
\emph{IEEE Access}, vol.~9, pp.~70961--71012, 2021.

\bibitem{zhou2024auth}
X.~Zhou, S.~Wang, K.~Wen, B.~Hu, X.~Tan, and Q.~Xie,
``Security-Enhanced Lightweight and Anonymity-Preserving User Authentication
Scheme for IoT-Based Healthcare,''
\emph{IEEE Internet of Things Journal}, vol.~11, no.~6, pp.~9718--9731, 2024.

\bibitem{kang2017iae}
B.~Kang, D.~Kim, and H.~Choo,
``Internet of Everything: A Large-Scale Autonomic IoT Gateway,''
\emph{IEEE Trans.\ Multi-Scale Computing Syst.}, vol.~3, pp.~206--214, 2017.

\bibitem{lajnef2022wsn}
N.~Lajnef, A.~Mami, and F.~Derbel,
``Applications of Wireless Sensor Networks and IoT Frameworks in Industry~4.0:
A Systematic Literature Review,''
\emph{Sensors}, vol.~22, no.~6, 2022.


\bibitem{alnahari2021auth}
W.~Alnahari and M.~T.~Quasim,
``Authentication of IoT Device and IoT Server using Security Key,''
in \emph{Proc.\ ICOTEN}, pp.~1--9, IEEE, 2021.

\bibitem{khatua2024puf}
M.~Khatua, S.~Das, and P.~M.~Rana,
``PUF-based Lightweight Mutual Authentication Scheme for IoT-based Smart
Grids,''
in \emph{Proc.\ IEEE ANTS}, pp.~1--6, IEEE, 2024.

\bibitem{sheik2023auth}
S.~A.~Sheik and A.~P.~Muniyandi,
``Secure Authentication Schemes in Cloud Computing with Glimpse of Artificial
Neural Networks: A Review,''
\emph{Cyber Security and Applications}, vol.~1, p.~100002, 2023.

\bibitem{he2023lightweight}
D.~He, Y.~Cai, S.~Zhu, Z.~Zhao, S.~Chan, and M.~Guizani,
``A Lightweight Authentication and Key Exchange Protocol with Anonymity for
IoT,''
\emph{IEEE Trans.\ Wireless Commun.}, vol.~22, no.~11, pp.~7862--7872, 2023.

\bibitem{roy2025design}
S.~Roy, M.~H.~Mahalat, and B.~Sen,
``Design and Implementation of Cost-Effective End-to-End Authentication
Protocol for PUF-Enabled IoT Devices,''
\emph{IEEE Trans.\ Emerging Topics in Computing}, 2025.

\bibitem{rosa2023lightsae}
P.~Rosa, A.~Souto, and J.~Cec\'{\i}lio,
``Light-SAE: A Lightweight Authentication Protocol for Large-Scale IoT
Environments Made with Constrained Devices,''
\emph{IEEE Trans.\ Network and Service Management}, vol.~20, no.~3,
pp.~2428--2441, 2023.

\bibitem{alruwaili2025raaf}
O.~Alruwaili, M.~Tanveer, S.~A.~Aldossari, S.~Alanazi, and A.~Armghan,
``RAAF-MEC: Reliable and Anonymous Authentication Framework for IoT-enabled
Mobile Edge Computing Environment,''
\emph{Internet of Things}, vol.~29, p.~101459, 2025.


\bibitem{yang2024sakms}
W.~Yang, C.~Hou, Y.~Wang, Z.~Zhang, X.~Wang, and Y.~Cao,
``SAKMS: A Secure Authentication and Key Management Scheme for IETF~6TiSCH
Industrial Wireless Networks based on Improved Elliptic-Curve Cryptography,''
\emph{IEEE Trans.\ Network Sci.\ Eng.}, vol.~11, no.~3, pp.~3174--3188, 2024.

\bibitem{li2020puf}
S.~Li, T.~Zhang, B.~Yu, and K.~He,
``A Provably Secure and Practical PUF-based End-to-End Mutual Authentication
and Key Exchange Protocol for IoT,''
\emph{IEEE Sensors Journal}, vol.~21, no.~4, pp.~5487--5501, 2020.

\bibitem{zheng2022puf}
Y.~Zheng, W.~Liu, C.~Gu, and C.-H.~Chang,
``PUF-based Mutual Authentication and Key Exchange Protocol for Peer-to-Peer
IoT Applications,''
\emph{IEEE Trans.\ Dependable and Secure Computing}, vol.~20, no.~4,
pp.~3299--3316, 2022.


\bibitem{alriyami2003cert}
S.~Al-Riyami and K.~G.~Paterson,
``Certificateless Public Key Cryptography,''
in \emph{Proc.\ ASIACRYPT}, vol.~2894, LNCS, pp.~452--473, Springer, 2003.

\bibitem{zhang2023blockchain}
P.~Zhang, Y.~Wang, G.~S.~Aujla, A.~Jindal, and Y.~D.~Al-Otaibi,
``A Blockchain-Based Authentication Scheme and Secure Architecture for
IoT-Enabled Maritime Transportation Systems,''
\emph{IEEE Trans.\ Intelligent Transportation Syst.}, vol.~24, no.~2,
pp.~2322--2331, 2023.

\bibitem{ahmed2023chains}
M.~T.~Al Ahmed, F.~Hashim, S.~J.~Hashim, and A.~Abdullah,
``Authentication-Chains: Blockchain-Inspired Lightweight Authentication
Protocol for IoT Networks,''
\emph{Electronics}, vol.~12, no.~4, p.~867, 2023.

\bibitem{sikarwar2020eff}
H.~Sikarwar, D.~Das, and S.~Kalra,
``Efficient Authentication Scheme Using Blockchain in IoT Devices,''
in \emph{Proc.\ AINA}, vol.~1151, pp.~630--641, Springer, 2020.


\bibitem{weerasena2024acc}
H.~Weerasena and P.~Mishra,
``Lightweight Multicast Authentication in NoC-based SoCs,''
in \emph{Proc.\ 25th ISQED}, pp.~1--8, IEEE, 2024.

\bibitem{blanchet2018proverif}
B.~Blanchet, B.~Smyth, V.~Cheval, and M.~Sylvestre,
``ProVerif 2.00: Automatic Cryptographic Protocol Verifier, User Manual and
Tutorial,''
Tech.\ Rep., INRIA, 2018.

\bibitem{dolev1983security}
D.~Dolev and A.~Yao,
``On the Security of Public Key Protocols,''
\emph{IEEE Trans.\ Information Theory}, vol.~29, no.~2, pp.~198--208, 1983.

\end{thebibliography}
\fi

\end{document}
        